% concatenated from Obstruction_full_pluri.tex
\documentclass{amsart}



%\usepackage[initials]{amsrefs}
\usepackage{amsmath,amsthm,amscd,amssymb}
\usepackage{mathrsfs}
\usepackage{graphicx}
%\usepackage[color]{showkeys}\definecolor{refkey}{rgb}{0.9451,0.2706,0.4941}\definecolor{labelkey}{rgb}{0.9451,0.2706,0.4941}

\usepackage{hyperref}
\usepackage{xcolor}
\hypersetup{
colorlinks=green,
citebordercolor=green,
linkbordercolor=green,
urlbordercolor=green
}


\usepackage{enumerate}
\usepackage{yhmath}%large widehat
\usepackage{color}
%%%%%%%%%%%%%%%%%%%%%%%%%%%%%%%%%%%%%
\usepackage{lineno}
%\pagewiselinenumbers
%\linenumbers
%%%%%%%%%%%%%%%%%%%%%%%%%%%%%%%%%%%%%


\theoremstyle{plain}
\newtheorem{thm}{Theorem}[section]
\newtheorem{cor}[thm]{Corollary}
\newtheorem{lem}[thm]{Lemma}
\newtheorem{prop}[thm]{Proposition}
\newtheorem{conj}[thm]{Conjecture}
\newtheorem{problem}[thm]{Problem}


\newtheorem{mtheorem}{Theorem}
\renewcommand{\themtheorem}{\Alph{mtheorem}}
\newtheorem{mcorollary}[mtheorem]{Corollary}


\theoremstyle{definition}
\newtheorem{dfn}[thm]{Definition}
\newtheorem{example}[thm]{Example}
\newtheorem*{idea}{Idea}
\newtheorem{question}[thm]{Question}
\theoremstyle{remark}
\newtheorem{remark}[thm]{Remark}
\newtheorem*{pf}{Proof}
\numberwithin{equation}{section}
%\numberwithin{figure}{section}
\usepackage{bm}

\newcounter{tmp}

\newcommand{\Ci}{\mathscr{C}}
\newcommand{\Sone}{\mathbb{S}^1}
\newcommand{\N}{\mathbb{N}}
\newcommand{\R}{\mathbb{R}}
\newcommand{\Z}{\mathbb{Z}}
\newcommand{\Li}{\mathcal{L}}
\newcommand{\lp}{\Lambda^{\prime}}
\newcommand{\la}{\Lambda}
\newcommand{\A}{A}
\newcommand{\W}{W}
\newcommand{\D}{D}
%\newcommand{\U}{\mathcal{U}}
\newcommand{\Op}{\mathcal{O}}
\newcommand{\di}{\displaystyle }


\newcommand{\esssup}{\operatorname*{ess\,sup}}
%\newcommand{\ubar}{\underaccent{\bar}}
\newcommand{\ubar}{\underline}


\def\smfd{W^{\mathrm{s}}}
\def\umfd{W^{\mathrm{u}}}
\def\lsmfd{W^{\mathrm{s}}_{\mathrm{loc}}}
\def\lssmfd{W^{\mathrm{ss}}_{\mathrm{loc}}}
\def\lumfd{W^{\mathrm{u}}_{\mathrm{loc}}}
\def\loc{\mathrm{loc}}
\def\Diff{\mathrm{Diff}}


%\def\setl{\setlength{\leftskip}{-18pt}}

\def\bb{\mathbb{B}}
\def\vv{\mathbb{V}}
\def\xx{\mathbb{X}}
\def\yy{\mathbb{Y}}
\def\gg{\mathbb{G}}
\def\uu{\mathbb{U}}
\def\bx{{\boldsymbol{x}}}
\def\by{{\boldsymbol{y}}}
\def\bv{{\boldsymbol{v}}}




\title[A robust obstruction to full strong pluripotency]{A robust obstruction to full strong pluripotency for wild blender-horseshoes}

\date{\today}


\author{Shin Kiriki}
\address[Shin Kiriki]{Department of Mathematics, Tokai 
University, 4-1-1 Kitakaname, Hiratsuka, Kanagawa, 259-1292, JAPAN}
\email{kiriki@tokai.ac.jp}

\author{Xiaolong Li}
\address[Xiaolong Li]{School of Mathematics and Statistics, 
Huazhong University of Science and Technology, Luoyu Road 1037, Wuhan, 430074, CHINA}
\email{lixl@hust.edu.cn}

\author{Yushi Nakano}
\address[Yushi Nakano]{
Department of Mathematics, Hokkaido University, Kita 10, Nishi 8, Kita-Ku, Sapporo, Hokkaido, 060-0810, 
JAPAN}
\email{yushi.nakano@math.sci.hokudai.ac.jp}

\author{Teruhiko Soma}
\address[Teruhiko Soma]{Department of Mathematical Sciences, 
Tokyo Metropolitan University, 1-1 Minami-Osawa, Hachioji, Tokyo, 192-0397, JAPAN}
\email{tsoma@tmu.ac.jp}



\subjclass[2020]{Primary: 37C20, 37C29; Secondary: 37C25}

\keywords{Takens' Last Problem, strong pluripotency, empirical measure, blender-horseshoe, homoclinic tangency}


\begin{document}



\begin{abstract}
Suppose that $M$ is a closed manifold of dimension greater than two and $r\geq 2$.
We show that there exists a $C^r$-diffeomorphism $f:M\longrightarrow M$ with a wild affine blender-horseshoe $\Lambda_f$ which is 
$C^r$-robustly and strongly pluripotent for $\Lambda_f^{(\mathrm{mj})}$ but not for $\Lambda_f$, where $\Lambda_f^{(\mathrm{mj})}$ is the subset of $\Lambda_f$ consisting of  elements with the majority condition.
Thus, within the present affine blender-horseshoe family in \cite{KNS}, there is a robust obstruction to extending the strong pluripotency from $\Lambda_f^{(\mathrm{mj})}$ to the whole horseshoe.
\end{abstract}


\maketitle



\section{Introduction}
Pluripotency is a term widely used in physiology and related fields to describe the ability of a system to transition from an undifferentiated state to multiple differentiated states governed by its internal dynamics.
In \cite{KNS}, this notion is formulated from the viewpoint of dynamical systems.
Specifically, it realizes some orbits initiating from a hyperbolic invariant set $\Lambda$ 
stochastically and observably under small perturbations of the dynamics.
We say that the system has full pluripotency if the dynamics realizes \emph{all} such orbits.
In this paper, we are interested in the subject whether given dynamics have full pluripotency or not.

Suppose that $M$ is a closed Riemannian manifold and $\mathrm{Diff}^r(M)$ $(r\geq 1)$ is the space of 
$C^r$-diffeomorphisms on $M$ with $C^r$-topology.
For any $f\in \mathrm{Diff}^r(M)$ and $x\in M$, consider the sequence of empirical  measures 
\begin{equation}\label{eqn_empirical}
\delta^{n}_{x,f}=
\frac{1}{n}\sum_{i=0}^{n-1}\delta_{f^i(x)},
\end{equation}
where $\delta_{f^i(x)}$ is the Dirac measure on $M$ supported at $f^i(x)$.
Let $\mathcal{P}(M)$ be the space of Borel probability measures on $M$ with weak$^*$-topology.
The notion of historic behavior is  introduced by Ruelle \cite{R01}.
We say that, for  $x\in M$, the $f$-forward orbit of $(f^n(x))_{n\geq 0}$ has 
\emph{historic behavior} if the sequence of empirical measures \eqref{eqn_empirical} does not converge in $\mathcal{P}(M)$.
A subset of $M$ is called a \emph{historic set} or a \emph{set with historic behavior} if it consists of initial points which give rise to orbits with historic behavior.
Note that the maximal historic set is Lebesgue measurable, see \cite[Proposition A.1]{KLNS22}.



The following problem proposed by Takens \cite{T08} is the main motivation of 
our study, which is nowadays called Takens' Last Problem 
(abbreviated as TLP for short).

\medskip
\emph{Are there persistent classes of smooth dynamical systems such that the maximal historic 
set has positive Lebesgue measure?}
\medskip

Here a non-empty subset $\mathcal{C}$ of $\mathrm{Diff}^r(M)$ is called a \emph{class}.
A property $\mathbf{P}$ on elements of $\mathrm{Diff}^r(M)$ is $\mathcal{C}$-\emph{persistent} or \emph{persistent with respect to} 
$\mathcal{C}$ if any elements of $\mathcal{C}$ satisfy the property.
Moreover, if the class $\mathcal{C}$ has non-empty interior, then $\mathbf{P}$ 
is called a $C^r$-\emph{robust} property of $\mathrm{Diff}^r(M)$.
TLP asks the existence of an observable historic set, 
where a subset of $M$ is called \emph{observable} if it has positive Lebesgue measure. 
An example of 2-dimensional dynamics with observable historic set is given by 
Colli-Vargas \cite{CV01}.
Their model is important in the sense that it provides a template for constructing 
diffeomorphisms with observable historic sets of special type, i.e.\ historic contracting wandering domains.
Here a non-empty connected open subset $D$ of $M$ is called a  \emph{wandering domain} 
if, for any $m,n\in \mathbb{N}\cup \{0\}$ $(m\neq n)$, 
$f^m(D)\cap f^n(D)=\emptyset$.
In fact, we present in \cite{KNS23} a class $\mathcal{C}$ of diffeomorphisms of dimension three or more with wild affine blender-horseshoes such that each element of $\mathcal{C}$ has a 
contracting wandering domain with historic behavior.
Here a blender-horseshoe $\Lambda$ is called \emph{wild} 
if $\umfd(\Lambda)$ and $\smfd(\Lambda)$ have 
a homoclinic tangency in the closure of the superposition region of $\Lambda$, 
see \cite[Section 6.2]{BDV05} for such a region.
By incorporating the Colli-Vargas model into any 2-dimensional diffeomorphism in the 
dissipative Newhouse domain $\mathcal{N}$ in $\mathrm{Diff}^r(M)$, it is shown 
in \cite{KS17} that there exists a dense subset $\mathcal{C}$ of $\mathcal{N}$ 
each element of which admits a contracting wandering domain with historic behavior.
This is an affirmative answer to TLP relative to a dense subset of $\mathcal{N}$.
Subsequently, other answers to TLP relative to certain persistent classes 
are given by Labouriau-Rodrigues \cite{LR17}, Barrientos \cite{B22}, Berge-Biebler \cite{BB23}, 
Jamilov-Mukhamedov \cite{JM} and so on.
 


Taking one more step beyond TLP, we will study detailed behavior of orbits from statistical or more strongly  geometrical points of view. 
Arguments on the Colli-Vargas model $f_0$ in \cite{CV01} implicitly imply that the geometrical average of orbits 
initiating any element of the whole horseshoe $\Lambda_{f_0}$ is realized observably by orbits of a diffeomorphism arbitrarily $C^r$-close to $f_0$, that is, 
$f_0$ has full strong pluripotency.
Moreover, it is shown in \cite{KLNSV} that any diffeomorphisms $f$ sufficiently close to 
$f_0$ have the the same property for the continuation $\Lambda_f$ of $\Lambda_{f_0}$.
This means that the strong full pluripotency is a $C^r$-robust property in 2-dimensional dynamics.
In the three or more dimensional case, the model given in \cite{KNS23} has 
strong pluripotency for the subset $\Lambda_f^{(\mathrm{mj})}$ of $\Lambda_f$ 
consisting of elements with the majority condition, 
see \eqref{eqn_mj_cond} in the next section for the definition.
Furthermore it is shown in \cite{KNS} that the strong pluripotency for $\Lambda_f^{(\mathrm{mj})}$ is $C^r$-robust.
The aim of this paper is to present three or more dimensional dynamics which are 
$C^r$-robustly and strongly pluripotent for $\Lambda_f^{(\mathrm{mj})}$ but not for $\Lambda_f$.








\section{Pluripotency and main theorem}

As in Introduction, suppose that $M$ is a closed Riemannian manifold and $r\geq 1$.
We consider the \emph{first Wasserstein metric} $d_W$ on the space $\mathcal{P}(M)$ of Borel probability measures on $M$ defined as
\[
d_W(\mu ,\nu )= \sup_\varphi \left\vert \int _M \varphi \, d\mu - \int _M \varphi \, d\nu \right\vert
\]
for $\mu$, $\nu\in \mathcal{P}(M)$, 
where the supremum is taken for all Lipschitz functions $\varphi : M \longrightarrow [-1, 1]$ 
with Lipschitz constant bounded by $1$. 
Note that the metric $d_W$ is compatible with the weak$^*$-topology on $\mathcal{P}(M)$.

Now we recall the concepts of pluripotency and strong pluripotency introduced in \cite{KNS}.


\begin{dfn}[Pluripotency]\label{dfn1}
Let $\la_f $ be a uniformly hyperbolic compact invariant set for $f \in \Diff^{r}(M)$.
\begin{enumerate}[(1)] 
\item $f$  is \emph{pluripotent} for  a subset $\la_f'$ of $\la_f$ if, for any $x\in \la_f'$,  
there exist $g \in \Diff^r(M)$ arbitrarily $C^r$-close to $f$ and 
a positive Lebesgue measure subset $\D_g$ of $M$ such that, 
for any $y\in \D_g$, 
\begin{equation}\label{defpl}
\lim _{n\to \infty} d_W( \delta_{y,g}^n, \delta_{x_g,g}^n ) =0,
\end{equation}
where $x_g\in \la_g'$ is the continuation of $x\in \la_f'$. 
\item
$f$ is \emph{strongly pluripotent} for $\la_f'$ if the 
following condition holds instead of \eqref{defpl} for $g$ and $D_g$ as above.
\begin{equation}\label{defspl}
\lim _{n\to \infty} 
\frac{1}{n}\sum_{i=0}^{n-1}
\sup_{y\in \D_{g}}\{\mathrm{dist}(g^i(y),g^i(x_{g}))\}=0.
\end{equation}
\end{enumerate}
Moreover such an $f$ is said to be \emph{fully} (strongly) pluripotent if $\la_f'$ is equal to the whole 
invariant set $\la_f$.
\end{dfn}
We note that \eqref{defspl} implies \eqref{defpl}, while the converse is not true in 
general.
For example, see \cite[Theorem 1.8]{KNS}.




In this paper we consider the case of $\dim M\geq 3$.
In \cite[Theorem B]{KNS}, we showed that 
there exists a neighborhood $\mathcal{U}_0$ of a diffeomorphism $f_0$ in $\mathrm{Diff}^r(M)$ with a wild affine blender-horseshoe 
such that any element $f$ of $\mathcal{U}_0$ has the continuation $\Lambda_f$ of 
$\Lambda_{f_0}$ and is strongly pluripotent for 
the subset $\Lambda_f^{(\rm mj)}$ of 
$\Lambda_f$ consisting of elements with the majority condition.
Here we say that an element $x$ of $\Lambda_f$ satisfies the \emph{majority condition} if 
\begin{equation}\label{eqn_mj_cond}
\liminf_{n\to\infty}\frac{
\#\left\{
j\ ;\ 
n- (3n)^{2/3} \leq j\leq n,\ 
v_{j}=0
\right\}
}{(3n)^{2/3}}\geq \frac12,
\end{equation}
where 
$(v_j)_{j\in \mathbb{Z}}\in \{0,1\}^{\mathbb{Z}}$ is the binary code corresponding to $x$.
Since $\Lambda_f^{({\rm mj})}$ is dense in but not equal to $\Lambda_f$, the following question  is natural.
%proposed in \cite[Question 1.10]{KNS} is natural.

\medskip
\emph{
Does there exist a diffeomorphism $f$ arbitrarily $C^r$-close to $f_0$ 
as  above and strongly pluripotent for the whole $\Lambda_{f}$?}
\medskip


The following theorem gives a negative answer to this question.


\begin{mtheorem}\label{thm_A}
Let $M$ be a closed manifold of dimension greater than two. 
Then there exist an element $f_0$ of $\mathrm{Diff}^r(M)$ $(r\geq 1)$ with an affine blender-horseshoe $\Lambda_{f_0}$ and a $C^r$-neighborhood $\mathcal{U}_0$ of $f_0$ satisfying the following condition. 
Any element $f$ of  $\mathcal{U}_0$ is not strongly pluripotent 
for the continuation $\Lambda_f$ of the whole $\Lambda_{f_0}$.
\end{mtheorem}

The diffeomorphism $f_0$ given in Theorem \ref{thm_A} is essentially the same as the one in \cite[Theorem~B]{KNS}, with the extra assumptions in Section \ref{S_2} added.
Theorem \ref{thm_A} exhibits an obstruction to full strong pluripotency in the present family.



By Theorem \ref{thm_A} together with \cite[Theorem B]{KNS}, we have the following corollary.

\begin{mcorollary}\label{cor_B}
If $2\leq r<\infty$, then there exist $f_0$ and $\mathcal{U}_0$ as in Theorem \ref{thm_A} 
such that  
any element $f$ of $\mathcal{U}_0$ is strongly pluripotent for $\Lambda_f^{({\rm mj})}$ but not 
for $\Lambda_f$.
\end{mcorollary}

We note that, to apply Theorem B in \cite{KNS}, $f_0$ in Corollary \ref{cor_B} must have a wild blender-horseshoe.

Corollary \ref{cor_B} raises the following question.


\begin{question}
Does there exist a maximal subset $\Lambda_{f_0}'$ of $\Lambda_{f_0}$ such that
every $f \in \mathcal{U}_0$ is strongly pluripotent for the continuation
$\Lambda_f'$ of $\Lambda_{f_0}'$?
If so, determine $\Lambda_{f_0}'$.
\end{question}



Theorem \ref{thm_A} says nothing about the (weak) pluripotency for $\Lambda_f$. 
So the following question remains open.


\begin{question}
Does there exist a diffeomorphism $f$ which is arbitrarily close to $f_0$ given in Theorem \ref{thm_A} 
and pluripotent for the whole $\Lambda_f$?
\end{question}



\section{Describability and Pluripotency Lemma}\label{S_describable}

In this section, we present the practical necessary and sufficient condition for a diffeomorphism $f$ on 
$M$ to be strongly pluripotent.

A pair $\{\mathbb{U}_0,\mathbb{U}_1\}$ of connected open sets in $M$ with mutually disjoint closures 
 is called a  \emph{coding pair} of a horseshoe $\Lambda_f$ of $f$ if 
\[
\Lambda_f =\bigcap_{i\in \mathbb{Z}} f^{i}(\mathbb U_0 \sqcup \mathbb U_1)
\]
and the restriction $f|_{\Lambda_f}:\Lambda_f\longrightarrow \Lambda_f$ is topologically conjugate 
to the shift map on $ \{ 0,1\} ^{\mathbb Z}$ by the coding map  
$\mathcal{I}_f: \Lambda_f  \longrightarrow  \{ 0,1\} ^{\mathbb Z}$ defined as 
\[
\left(\mathcal{I}_f(x) \right)_j = v \quad \text{if $f^j(x) \in \mathbb U_{v}$}, 
\]
where $(\mathcal{I}_f(x) )_j$ stands for the $j$-th entry of $\mathcal{I}_f(x)$. 




\begin{dfn}[Describability]\label{describable}
Let $\Sigma$ be a subset of $\{0,1\}^{\mathbb N_0}$ and 
$f$ an element of $\mathrm{Diff}^{r}(M)$ with a horseshoe $\Lambda_f$ associated with a coding pair $\{ \mathbb U_0, \mathbb U_1 \}$, 
where $\mathbb N_0=\mathbb{N}\cup \{0\}$.  
We say that $f$ is $\Sigma$-\emph{describable} over $\Lambda_f$  
if any element $\underline{v}=(v_{0}v_{1}v_{2}\ldots\,)$ of $\Sigma$
satisfies the following conditions: 

\begin{enumerate}[({D}1)]
\makeatletter
\renewcommand{\p@enumi}{D}
\makeatother
%\setl
\item\label{D1}
There exist a strictly increasing sequence $(\alpha _k)_{k\in \mathbb N}$ of non-negative integers and 
integers $\beta _k$ $(k\in \mathbb N)$ with $0\leq \beta _k\leq \alpha _{k+1}-\alpha _k-1$ and  
such that 
\[
\lim _{n\rightarrow \infty} \frac{\#\left\{j\,;\, 0\le j \le n-1,\ j \in \bigcup _{k=1}^\infty \mathbb{I}_k\right\}}{n} =1, 
\]
where $\mathbb{I}_k
=[\alpha_k, \alpha _k+ \beta_k] \cap \mathbb Z$.
\item\label{D2}
There exist an element $g$ of $\mathrm{Diff}^{r}(M)$ arbitrarily $C^r$-close to $f$ and  
a positive Lebesgue measure subset $D_g$ of $M$  
such that  
\[
g^j (D_g) \subset \mathbb U_{v_{j}}
\]
for any  $j\in \bigcup_{k\in \mathbb{N}
} \mathbb I_k$.  
\end{enumerate}
\end{dfn}




The following theorem  in \cite[Theorem A]{KNS} plays 
an important role in the proof of the main theorem, 
where $\widehat \Sigma$ is the subset of $\{0,1\}^{\mathbb{Z}}$ consisting of elements 
$(v_j)_{j\in \mathbb{Z}}$ with $(v_j)_{j\in \mathbb{N}_0}\in \Sigma$.


\begin{thm}[Pluripotency Lemma]\label{p-lemma}
Suppose that $f$ is an element of $\mathrm{Diff}^r(M)$ $(r\geq 1)$ with a horseshoe $\Lambda_f$ associated with a coding pair $\{ \mathbb U_0, \mathbb U_1 \}$ and 
$\Sigma$ is a non-empty subset of $\{0,1\}^{\mathbb N_0}$.
Then $f$ is $\Sigma$-describable if and only if 
$f$ is strongly pluripotent for $\mathcal{I}_f^{-1}(\widehat \Sigma)$.
\end{thm}





\begin{remark}
The novelty of Theorem \ref{thm_A} lies in using the converse direction of the Pluripotency Lemma as an obstruction principle. 
In our previous works, one verified describability in order to prove strong pluripotency. 
In contrast, here we show that, for asymptotically 1-filling itineraries (see \eqref{eqn_1-filling} for the difinition), describability would force long good blocks whose projected-area growth is incompatible with 
the geometric phenomena caused by the $g$-orbit of $D_g$. 
This obstruction rules out full strong pluripotency.
\end{remark}



\section{Wild affine blender-horseshoes}\label{S_2}

For simplicity, we only consider the case that $M$ is a closed manifold of dimension three.
In the case of $\dim M=n>3$, our arguments still work    
for certain elements of $\mathrm{Diff}^r(M)$ having a blender-horseshoe 
$\Lambda$ with $\dim W^{\mathrm{u}}(\Lambda)=\dim W^{\mathrm{cs}}(\Lambda)=1$ and $\dim W^{\mathrm{ss}}(\Lambda)=n-2$.
A diffeomorphism $f_0$ on $M$ presented in this section has a wild affine blender-horseshoe similar to those in \cite{KNS23, KNS}.

\medskip

 
Let $\lambda_{\rm ss}, \lambda_{\rm cs0}, \lambda_{\rm cs1}$ and $\lambda_{\rm u}$ 
be real positive constants with
\begin{equation}\label{eqn_eigen_v}
\lambda_{\rm ss}<\frac12,\quad
\lambda_{\rm ss}<\lambda_{\rm cs0}\leq \lambda_{\rm cs1}<1<\lambda_{\rm cs0}+\lambda_{\rm cs1},\quad 2<\lambda_{\rm u}.
\end{equation}
Here the inequalities $\lambda_{\rm ss}<2^{-1}$ and $2<\lambda_{\rm u}$ are necessary to define a blender-horseshoe.
The condition $1<\lambda_{{\rm cs}0}+\lambda_{{\rm cs}1}$ is used to show that the blender-horseshoe has the \emph{distinctive property}.  
See Section 6.2 in \cite{BDV05} for details.
The property is required in the proof of \cite[Theorem B]{KNS}, 
which is essential for proving Corollary \ref{cor_B}.
It follows from \eqref{eqn_eigen_v} that   
\begin{equation}\label{eqn_lamcs1}
\lambda_{\rm cs1}\lambda_{\rm u}>1,
\end{equation}
which is crucial in the proof of Theorem \ref{thm_A}.
We fix a sufficiently small  $\varepsilon_0>0$.
For example, it is taken so that $1/2-\lambda_{\rm ss}(1+\varepsilon_0)$ is positive.
The closed interval $[-\varepsilon_{0},1+\varepsilon_{0}]$ is denoted by  $I_{\varepsilon_{0}}$ 
and the cube $I_{\varepsilon_{0}}^{3}=I_{\varepsilon_{0}}\times I_{\varepsilon_{0}}\times I_{\varepsilon_{0}}$ by $\mathbb{B}$.
We suppose that $\mathbb{B}$ is embedded in $M$ 
and $M$ has a Riemannian metric extending the standard Euclidean metric on $\mathbb{B}$.
Consider the sub-blocks of $\mathbb{B}$ defined as 
\[
\mathbb{V}_{0}= [-\varepsilon_{0}, \lambda_{\rm u}^{-1}+\varepsilon_0]\times I_{\varepsilon_{0}}^{2},\quad 
\mathbb{V}_{1}= [1-\lambda_{\rm u}^{-1}-\varepsilon_0,1+\varepsilon_{0}]\times I_{\varepsilon_{0}}^{2}. 
\]
Let $f_{0}$ be a $C^r$-diffeomorphism on $M$ such that $f_{0}|_{\mathbb{V}_{0}\cup \mathbb{V}_{1}}$ is defined as
\begin{equation}\label{eqn_f0V}
f_{0}(x,y,z)=
\begin{cases}
(\lambda_{\mathrm{u}}x,\lambda_{\mathrm{ss}}y,\zeta_0(z))&\text{if}\ (x,y,z)\in \mathbb{V}_{0},\\
(\lambda_{\mathrm{u}}(1-x),1-\lambda_{\mathrm{ss}}y,\zeta_1(z)) &\text{if}\ (x,y,z)\in \mathbb{V}_{1},
\end{cases}
\end{equation} 
where 
$\zeta_0$ and $\zeta_1$ are the affine maps on $I_{\varepsilon_0}$ given by
%\begin{equation}\label{eqn_def_zeta}
\[
\zeta_0(z)=\lambda_{\mathrm{cs} 0}z\quad\text{and}\quad \zeta_1(z)=\lambda_{\mathrm{cs} 1}z+1-\lambda_{\mathrm{cs} 1}.
\]
%\end{equation}
From our setting, $f_{0}$ has the uniformly hyperbolic set  
\[
\Lambda_{f_{0}}=\bigcap_{n\in \mathbb{Z}} f_{0}^n(\mathbb{V}_{0}\cup \mathbb{V}_{1}),
\]
which belongs to the class of affine blender-horseshoes, see \cite{BD96,BD12} for details. 
Note that $p_{f_0}=(0,0,0)$ is the saddle fixed point of $f_0$ contained in $\mathbb{V}_0$.







Let $a_{1},a_{2},a_{3},a_{4}$ and $\mu$ be constants 
satisfying  
\begin{subequations}
\begin{align}
&a_{1}>0,\quad a_2>0,\quad a_3a_4<0,\label{eqn_aaa}\\
&|a_{3}|\left(\frac12+\varepsilon_0\right)<\frac12-\lambda_{\rm ss}(1+\varepsilon_0),\label{eqn_|a3|}\\
&|a_3|<\min\left\{\dfrac{a_2|a_4|}{\sqrt{4a_1^2+a_2^2}},\dfrac{a_1a_2}{\sqrt{4a_1^2+a_2^2}}(1-2\varepsilon_0)\right\},\label{eqn_min}\\
&a_1^{-1}a_2(1+\varepsilon_0+\mu)^2<\left(\frac12-\lambda_{\rm u}^{-1}\right)^2,\  
a_1^{-1}a_2a_4^2(1+\varepsilon_0+\mu)<\left(\frac12+\varepsilon_0\right)^2,\label{eqn_a42}\\
&0<-\varepsilon_0+\mu,\quad a_2(1+\varepsilon_0+\mu)<\lambda_{\rm u}^{-1}.\label{eqn_e_0mu}
\end{align}
\end{subequations}
It is not hard to find such constants by taking $a_1$ sufficiently large and next $|a_3|$ sufficiently small.
The conditions other than \eqref{eqn_min} are used in this section, but 
\eqref{eqn_min} is in the proof of Lemma \ref{l_ea_horizontal}.

Now we consider the subsets 
\begin{align*}
\mathbb{X}_{0}&= [0, \lambda_{\rm u}^{-1}]\times I_{\varepsilon_{0}}^{2},\quad 
\mathbb{X}_{1}= [1-\lambda_{\rm u}^{-1},1]\times I_{\varepsilon_{0}}^{2},\\
\mathbb{X}_2&=\left\{(x,y,z)\in \mathbb{B}\,\Big|\,z\geq \frac{a_1}{a_2}\left(x-\frac12\right)^2-\mu\right\}
\end{align*}
of $\mathbb{B}$.
Note that $\mathbb{X}_{0}\subset \mathbb{V}_{0}$ and $\mathbb{X}_{1}\subset \mathbb{V}_{1}$.


We suppose that the $y$-axis direction in $\mathbb{R}^3$ is vertical and 
say that a surface in $\mathbb{R}^3$ is parallel to the $xz$-plane is \emph{horizontal}.
Let $q_+$, $q_-$ be the intersection points of the parabolic cylinder 
$z=\dfrac{a_1}{a_2}\left(x-\dfrac12\right)^2-\mu$, the vertical plane $z=1+\varepsilon_0$ and 
the horizontal plane $y=\frac12$ in $\mathbb{R}^3$.
Since the $x$-entry of $q_\pm$ is $\dfrac12\pm\sqrt{a_1^{-1}a_2(1+\varepsilon_0+\mu)}$, 
the former inequality of \eqref{eqn_a42} implies $\xx_2\cap (\xx_0\cup \xx_1)=\emptyset$.
We denote by $\mathbb{Y}$ and $\mathbb{S}$ the closures of $[\lambda_{\rm u}^{-1}, 1-\lambda_{\rm u}^{-1}]\times I_{\varepsilon_0}^2\setminus \mathbb{X}_2$ and $\mathbb{B}\setminus [0,1]\times I_{\varepsilon_0}^2$ in $\mathbb{B}$ respectively.
The former inequality of \eqref{eqn_e_0mu} implies that $\mathbb{Y}$ consists of two components 
and hence the frontier of $\mathbb{X}_2$ in $\mathbb{B}$ consists of the two curved rectangles $F_2^{\pm}$ 
as illustrated in Figure \ref{f_X012}.
Here, for a subset $A$ of a topological space $B$,  the \emph{frontier} (or \emph{topological boundary}) of $A$ in $B$ is the intersection 
$\mathrm{Cl}(A)\cap \mathrm{Cl}(B\setminus A)$, which is denoted by $\mathrm{Fr}_B(A)$.
%%%%%%%%%%%%%%%%%%%%%%%%%%%%%%%%%%%%%%%%%%%%%%%%%%%%%%%
\begin{figure}[hbtp]
\centering
\scalebox{0.6}{\includegraphics[clip]{f_1}}
\caption{View from the top. $\mathrm{Fr}(\xx_2)=F_2^-\cup F_2^+$.}
\label{f_X012}
\end{figure}
%%%%%%%%%%%%%%%%%%%%%%%%%%%%%%%%%%%%%%%%%%%%%%%%%%%%%%%
Simply we set $\mathrm{Fr}_{\mathbb{B}}(A)=\mathrm{Fr}(A)$ for any subset $A$ of $\mathbb{B}$.
Then $\mathrm{Fr}(\xx_0)=\{0,\lambda_{\rm u}^{-1}\}\times I_{\varepsilon_0}^2$, 
$\mathrm{Fr}(\xx_1)=\{1-\lambda_{\rm u}^{-1},1\}\times I_{\varepsilon_0}^2$ and 
$\mathrm{Fr}(\mathbb{S})=\{0,1\}\times I_{\varepsilon_0}^2$, 


For any $(x,y,z)\in \mathbb{X}_{2}$ and constants $a_1$, $a_2$, $a_3$, $a_4$, $\mu$ satisfying  \eqref{eqn_aaa}--\eqref{eqn_e_0mu}, 
we suppose that   
\begin{equation}\label{eqn_tangency}
f_{0}^{2}(x,y,z)=\left(-a_{1}\Bigl(x-\frac{1}{2}\Bigr)^{2}+a_{2}(z+\mu),\ a_{3}\Bigl(y-\frac{1}{2}\Bigr)+\frac{1}{2},\ 
a_{4}\Bigl(x-\frac{1}{2}\Bigr)+\frac{1}{2} \right).
\end{equation}


The latter inequality of \eqref{eqn_a42} implies that $f_0^2(q_\pm)$ are contained in 
$\{0\}\times \mathrm{Int} I_{\varepsilon_0}^2$.
By this fact together with \eqref{eqn_e_0mu}, one can prove that $f_0^2(\mathbb{X}_2)$ is 
contained in $\mathbb{X}_0$.
See Figures \ref{f_blender} and \ref{f_Yf0} below.
The image of $f_0^2(\xx_2)$ by the orthogonal projection to the $y$-axis is the closed interval 
$\left[-|a_3|\left(\frac12+\varepsilon_0\right)+\frac12,\,  |a_3|\left(\frac12+\varepsilon_0\right)+\frac12\right]$.
Hence, by \eqref{eqn_|a3|}, $f_0^2(\xx_2)$ lies above $f_0(\xx_0)$ and below $f_0(\xx_1)$.
See Figure \ref{f_blender}. 
%%%%%%%%%%%%%%%%%%%%%%%%%%%%%%%%%%%%%%%%%%%%%%%%%%%%%%%
\begin{figure}[hbtp]
\centering
\scalebox{0.6}{\includegraphics[clip]{f_2}}
\caption{$h_\pm=\pm |a_3|\left(\frac12+\varepsilon_0\right)+\frac12$. 
The union of shaded rectangles represents $f_0^2(\mathrm{Fr}(\xx_2))$.
The front square $\{0\}\times I_{\varepsilon_0}^2$ of $\overline{\mathbb{B}\setminus \mathbb{S}}$  is the component of $W^{\mathrm{s}}(p_{f_0},f_0)\cap \mathbb{B}$ containing $p_{f_0}$.}
\label{f_blender}
\end{figure}
%%%%%%%%%%%%%%%%%%%%%%%%%%%%%%%%%%%%%%%%%%%%%%%%%%%%%%%
From our definition \eqref{eqn_f0V} of $f_0$ on $\mathbb{V}_0\cup \mathbb{V}_1$, we know that the $f_0^2$-image of the frontier $\mathrm{Fr}(\mathbb{X}_0\cup \mathbb{X}_1)$ is also contained in 
$\{0\}\times \mathrm{Int} I_{\varepsilon_0}^2$.
As in \cite[Lemma 1.4]{KNS23}, one can choose $\mu$ in \eqref{eqn_tangency} 
so that $f_0$ has a $C^{1}$-robust homoclinic tangency associated with $\Lambda_{f_0}$. 
Then $\Lambda_{f_0}$ is called a \emph{wild affine blender-horseshoe}.


By \eqref{eqn_tangency}, for $\boldsymbol{x}=(x,y,z)\in \mathbb{X}_2$, 
\begin{equation}\label{eqn_DfX2}
Df_0^2(\bx)=\begin{pmatrix}-2a_1\left(x-\frac12\right)&0&a_2\\
0&a_3&0\\
a_4&0&0\end{pmatrix}.
\end{equation}
By \eqref{eqn_aaa}, $\det Df_0^2(\bx)=-a_2a_3a_4>0$.


By \eqref{eqn_f0V}, for any horizontal surface $H$ in $\xx_i$ $(i=0,1)$ and a 
Lebesgue measurable subset $A$ of $H$, $f_0(H)$ is a horizontal surface in $\mathbb{B}$ and 
\begin{equation}\label{eqn_Leb_P0}
\mathrm{Leb}_{f_0(H)}(f_0(A))=\lambda_{{\rm cs}i}\lambda_{\rm u}\mathrm{Leb}_H(A),
\end{equation}
where $\mathrm{Leb}_H(\cdot)$ denotes the standard 2-dimensional Lebesgue measure on $H$.
On the other hand, by \eqref{eqn_tangency} and \eqref{eqn_DfX2}, for any horizontal surface $H'$ in $\xx_2$ and a Lebesgue measurable subset $A'$ of $H'$, 
$f_0^2(H')$ is a horizontal surface in $\xx_0$ and 
\begin{equation}\label{eqn_Leb_P}
\mathrm{Leb}_{f_0^2(H')}(f_0^2(A'))=|a_2a_4|\mathrm{Leb}_{H'}(A').
\end{equation}
 
 


We denote by $\mathrm{proj}:\mathbb{R}^3\longrightarrow \mathbb{R}^2$ the orthogonal projection to the $xz$-plane, that is, $\mathrm{proj}(x,y,z)=(x,z)$.
For any subset $K$ of $\mathbb{R}^3$, we set $\mathrm{proj}(K)=\widehat K$.

From now on, we work under the following extra assumptions.


\subsection*{Extra Assumptions}
The intersection $f_0(\xx_2)\cap \mathbb{B}$ is empty.
There exists an arc $\alpha$ in $M$ with $\mathrm{Int}\,\alpha\subset W^{\rm u}(p_{f_0},f_0)$ 
and connecting $p_{f_0}$ with an attracting fixed point  
$s_0$ of $f_0$ outside $\mathbb{B}$.
The sequence 
$f_0^n(\yy\cup \mathbb{S})$ converges to $\alpha$ as $n\to \infty$ with respect to the Hausdorff metric.
See Figure \ref{f_Yf0}. 
%%%%%%%%%%%%%%%%%%%%%%%%%%%%%%%%%%%%%%%%%%%%%%%%%%%%%%%
\begin{figure}[hbtp]
\centering
\scalebox{0.6}{\includegraphics[clip]{f_3}}
\caption{Schematic figure just presenting the connections of $f_0^2(\xx_i)$ $(i=0,1,2)$ and $f_0^2(\yy)$. 
They do not indicate their exact placements in $M$.
}
\label{f_Yf0}
\end{figure}
%%%%%%%%%%%%%%%%%%%%%%%%%%%%%%%%%%%%%%%%%%%%%%%%%%%%%%%



\section{Diffeomorphisms close to $f_0$}\label{S_near_f0}

\subsection{Extensions of $f_0|_{\mathbb{X}_0}$, $f_0|_{\mathbb{X}_1}$, 
$f_0^2|_{\mathbb{X}_2}$ to diffeomorphisms on $\mathbb{R}^3$}

By \eqref{eqn_f0V}, each component of $\mathrm{Fr}(\mathbb{X}_i)$ $(i=0,1)$ is a component of $\smfd(p_{f_0},f_0)\cap \mathbb{B}$.
Since $f_0^2(\mathrm{Fr}(\mathbb{X}_2))$ is contained in $\{0\}\times I_{\varepsilon_0}^2\subset \smfd(p_{f_0},f_0)$, $\mathrm{Fr}(\mathbb{X}_2)$ is also a union of two  
components of $\smfd(p_{f_0},f_0)\cap \mathbb{B}$.


Consider a small neighborhood $\mathcal{U}_0$ of $f_0$ in $\mathrm{Diff}^r(M)$.
For any $g\in \mathcal{U}_0$, we denote by $p_g$ the saddle fixed point of $g$ obtained as the continuation of $p_{f_0}$.
For $i=0,1,2$, let $\xx_{g,i}$  be the compact subset of $\mathbb{B}$ such that $\mathrm{Fr}(\xx_{g,i})\subset \smfd(p_g,g)\cap \mathbb{B}$ is the continuation of $\mathrm{Fr}(\xx_i)$.
The compact subset $\yy_g$ of $\mathbb{B}$ corresponding to $\yy$ 
is defined similarly.
We denote by $\mathbb{S}_g$ the closure of the union of components of $\mathbb{B}\setminus (\xx_{g,0}\cup \xx_{g,1})$ disjoint from $\xx_{g,2}$.
One can suppose that $\yy_g\cup \mathbb{S}_g$ still satisfies Extra Assumptions with respect to $g$ instead of $f_0$ 
if necessary by replacing $\mathcal{U}_0$ with a smaller neighborhood of $f_0$.
 



For $i=0,1$, we consider the diffeomorphism $\varphi_i:\mathbb{R}^3\longrightarrow \mathbb{R}^3$ which has the same form as $f_0|_{\vv_i}$ in \eqref{eqn_f0V} for \emph{any} $(x,y,z)\in \mathbb{R}^3$.
In particular, $\varphi_i|_{\xx_i}=f_0|_{\xx_i}$ holds.
Note that $\varphi_0$ is a linear map on $\mathbb{R}^3$ and $\varphi_1$ 
is an affine map fixing $(\lambda_{\mathrm u}(1+\lambda_{\mathrm u})^{-1},
(1+\lambda_{\mathrm s})^{-1},1)$, 
each of which preserves horizontal surfaces in $\mathbb{R}^3$.


Let $\eta:\mathbb{R}\longrightarrow \mathbb{R}$  be the unimodal $C^1$-function defined as 
\[
\eta(x)=\begin{cases}
2a_1\left(\dfrac12+\varepsilon_0\right)x+a_1\left(\dfrac12+\varepsilon_0\right)\!\left(-\dfrac12+\varepsilon_0\right) &
\text{if $x<-\varepsilon_0$},\\[4pt]
-a_1\left(x-\dfrac12\right)^2&\text{if $-\varepsilon_0\leq x\leq 1+\varepsilon_0$},\\[4pt]
-2a_1\left(\dfrac12+\varepsilon_0\right)x+a_1\left(\dfrac12+\varepsilon_0\right)\!\left(\dfrac32+\varepsilon_0\right) &
\text{if $x>1+\varepsilon_0$}.
\end{cases}
\]
From the definition, we have
\begin{equation}\label{eqn_d_eta}
\left|\frac{d\eta}{dx}(x)\right|\leq 2a_1\left(\frac12+\varepsilon_0\right)
\end{equation}
for  any $x\in\mathbb{R}$.
By using the function $\eta$, we define the $C^1$-diffeomorphism $\varphi_2:\mathbb{R}^3\longrightarrow \mathbb{R}^3$ by    
\begin{equation}\label{eqn_phi_2}
\varphi_{2}(x,y,z)=\left(\eta(x)+a_{2}(z+\mu),\ a_{3}\Bigl(y-\frac{1}{2}\Bigr)+\frac{1}{2},\ 
a_{4}\Bigl(x-\frac{1}{2}\Bigr)+\frac{1}{2} \right).
\end{equation}
By \eqref{eqn_tangency}, $\varphi_2|_{\xx_2}$ is equal to $f_0^2|_{\xx_2}$.
From \eqref{eqn_d_eta}, we know that the $(1,1)$ entry of $D\varphi_2(\bx)$ $(\bx\in \mathbb{R}^3)$ is  uniformly bounded unlike that of \eqref{eqn_DfX2}.
This is the reason why we take the form \eqref{eqn_phi_2} of $\varphi_2$ different from \eqref{eqn_tangency} of $f_0^2|_{\xx_2}$ on $\mathbb{R}^3\setminus \xx_2$. 
Note that $\varphi_2$ also preserves horizontal surfaces in $\mathbb{R}^3$.



\subsection{Almost horizontal surfaces in $\mathbb{R}^3$}
To prove Theorem \ref{thm_A}, for any $g \in \mathrm{Diff}^r(M)$ sufficiently $C^r$-close to $f_0$, one must examine every subset $D_g \subset M$ of positive Lebesgue measure satisfying condition \eqref{D2} in Definition \ref{describable}.
Since such subsets lack an explicit characterization, it is generally impossible to determine directly whether the 3-dimensional Lebesgue measure of $g^n(D_g)$ diverges as $n \to \infty$. 
To circumvent this difficulty, we instead consider a certain horizontal section $A$ of $D_g$ and demonstrate that its 2-dimensional Lebesgue measure diverges to infinity. 
As $n \to \infty$, the geometry of $g^n(A)$ typically becomes highly complex. 
To address this issue, we investigate a sequence of almost horizontal surfaces containing $g^n(A)$, rather than analyzing the iterates $g^n(A)$ directly.




Let $\varepsilon$ be a positive number sufficiently smaller than $\varepsilon_0$, which will be 
fixed in the proof of Lemma \ref{l_ea_horizontal}.
We consider the $\mathrm{uc}$-cone-field 
\[
\boldsymbol{C}_{\varepsilon}^{\mathrm{uc}}(\boldsymbol{x})=
\left\{
\boldsymbol{v}=(v^{\mathrm{u}},v^{\mathrm{s}},v^{\mathrm{cs}})\in T_{\boldsymbol{x}}(\mathbb{R}^3)\,;\,
|v^{\mathrm{s}}|\leq\varepsilon\sqrt{(v^{\mathrm{u}})^2+(v^{\mathrm{cs}})^2}
\right\}
%\end{align*}
\]
on $\mathbb{R}^3$.
A surface $S$ in $\mathbb{R}^3$ is said to be $\varepsilon$-\emph{almost horizontal} if, for any 
$\bx\in S$, $T_{\bx}S$ is contained in $\boldsymbol{C}_{\varepsilon}^{\mathrm{uc}}(\bx)$.
Furthermore such a surface $S$ is called \emph{proper} 
if the restriction $\mathrm{proj}|_S:S\longrightarrow \mathbb{R}^2$ is a proper map onto $\mathbb{R}^2$, that is, 
$(\mathrm{proj}|_S)^{-1}(K)$ is compact for any compact subset $K$ of $\mathbb{R}^2$.
This condition is necessary to avoid the situation where the image $\widehat A$ is multi-folded in $\mathbb{R}^2$ for a subset $A$ of $S$.
In fact, since $\mathrm{proj}|_S$ is locally homeomorphic for any $\varepsilon$-almost horizontal proper surface $S$,  it follows from the standard fact of 
covering theory that $\mathrm{proj}|_S:S\longrightarrow \mathbb{R}^2$ is a surjective homeomorphism.
The $\varepsilon$-almost horizontality of $S$ will be used to show the second inequality of \eqref{eqn_Leb_A} below.



Fix connected and mutually disjoint open neighborhoods $\mathbb{U}_i$ of 
$\xx_i$ $(i=0,1,2)$ in $(-\varepsilon_0,1+\varepsilon_0)\times (-2\varepsilon_0,1+2\varepsilon_0)^2$.
We may assume that, for any $g\in \mathcal{U}_0$, $\mathbb{U}_i$ contains $\xx_{g,i}$ and 
$\{\uu_0,\uu_1\}$ is a coding pair for $\Lambda_g$.
Moreover one can suppose that $g(\mathbb{U}_i)$ $(i=0,1)$ is contained in $\mathbb{B}$.
There exist diffeomorphisms $\varphi_{g,i}:\mathbb{R}^3\longrightarrow \mathbb{R}^3$ such that 
$\varphi_{g,i}|_{\xx_{g,i}}=g|_{\xx_{g,i}}$ for $i=0,1$, 
$\varphi_{g,2}|_{\xx_{g,2}}=g^2|_{\xx_{g,2}}$ 
and 
$\varphi_{g,j}|_{\mathbb{R}^3\setminus \mathbb{U}_j}=\varphi_j|_{\mathbb{R}^3\setminus \mathbb{U}_j}$ for $j=0,1,2$.
By taking $\mathcal{U}_0$ sufficiently small, we may assume that $\varphi_{g,i}$ $(i=0,1)$ is arbitrarily $C^r$-close to $\varphi_i$ and $\varphi_{g,2}$ is arbitrarily $C^1$-close to $\varphi_2$.



\begin{lem}\label{l_ea_horizontal}
There exists a neighborhood $\mathcal{U}_0$ of $f_0$ such that, 
for any $\varepsilon$-almost horizontal proper surface $S$ in $\mathbb{R}^3$, $g\in\mathcal{U}_0$ and $i=0,1,2$, 
the surface 
$\varphi_{g,i}(S)$ is also $\varepsilon$-almost horizontal and proper.
\end{lem}
\begin{proof}
First we consider the case of $i=0,1$.
By \eqref{eqn_f0V}, one can choose $\mathcal{U}_0$ so that 
\begin{equation}\label{eqn_Dphi_i}
D\varphi_{g,i}(\bx)=
\begin{pmatrix} (-1)^i\lambda_{\rm u}+\alpha_{11}(\bx)& \alpha_{12}(\bx)& \alpha_{13}(\bx)\\
\alpha_{21}(\bx)&(-1)^i \lambda_{\rm ss}+\alpha_{22}(\bx)& \alpha_{23}(\bx)\\
\alpha_{31}(\bx)& \alpha_{32}(\bx) &\lambda_{{\rm cs}i}+\alpha_{33}(\bx)
\end{pmatrix}
\end{equation}
holds for any $g\in\mathcal{U}_0$ and $\bx=(x,y,z)\in \mathbb{R}^3$,  
where each $\alpha_{kl}:\mathbb{R}^3\longrightarrow \mathbb{R}$ is a continuous function supported on $\mathbb{U}_i$ and satisfying  
$\sup\bigl\{|\alpha_{kl}(\bx)|\, ;\, \bx \in \mathbb{R}^3\bigr\}=O(\varepsilon^2)$.
Here we denote by $O(\varepsilon^2)$ a non-negative constant with $\limsup\limits_{\varepsilon\to +0}
O(\varepsilon^2)/\varepsilon^2<\infty$.
The definition of $\boldsymbol{C}_{\varepsilon}^{\mathrm{uc}}(\boldsymbol{x})$ 
implies $|v^{\mathrm{s}}|\leq \varepsilon$ for any $\bv=(v^{\mathrm{u}},v^{\mathrm{s}},v^{\mathrm{cs}})\in 
\boldsymbol{C}_{\varepsilon}^{\mathrm{uc}}(\boldsymbol{x})$ with 
$(v^{\mathrm{u}})^2+(v^{\mathrm{cs}})^2=1$.
We set $D\varphi_{g,i}(\bx)(v^{\rm u},v^{\rm s},v^{\rm cs})^T
=(\widehat{v}^{\rm u},\widehat{v}^{\rm s},\widehat{v}^{\rm cs})^T$, 
where $(v_1,v_2,v_3)^T$ denotes the transposed column vector of a row vector $(v_1,v_2,v_3)$.
By \eqref{eqn_Dphi_i}, 
\[
|\widehat{v}^{\rm s}|\leq \lambda_{\rm ss}|v^{\rm s}|+O(\varepsilon^2)
\leq \varepsilon\lambda_{\rm ss}+O(\varepsilon^2).
\]
Moreover, 
\begin{align*}
\varepsilon\sqrt{(\widehat{v}^{\rm u})^2+(\widehat{v}^{\rm cs})^2}
&\geq 
\varepsilon\sqrt{\lambda_{\rm u}^2(v^{\rm u})^2+\lambda_{{\rm cs}i}^2(v^{\rm cs})^2}-O(\varepsilon^2)\\
&>\varepsilon\lambda_{{\rm cs}i}\sqrt{(v^{\rm u})^2+(v^{\rm cs})^2}-O(\varepsilon^2)\\
&
=\varepsilon\lambda_{{\rm cs}i}-O(\varepsilon^2).
\end{align*}
By \eqref{eqn_eigen_v},   
$|\widehat{v}^{\rm s}|\leq \varepsilon\sqrt{(\widehat{v}^{\rm u})^2+(\widehat{v}^{\rm cs})^2}$ holds 
if we take $\varepsilon$ sufficiently small.
It follows that $D\varphi_{g,i}(\bx)\bv\in \boldsymbol{C}_{\varepsilon}^{\mathrm{uc}}(\boldsymbol{x})$.




In the case of $i=2$, by \eqref{eqn_phi_2}, we have 
\begin{equation}\label{eqn_Dphi2}
D\varphi_{g,2}(\bx)=\begin{pmatrix} \dfrac{d\eta}{dx}(x)+\beta_{11}(\bx)&\beta_{12}(\bx)&a_2+\beta_{13}(\bx)\\
\beta_{21}(\bx)&a_3+\beta_{22}(\bx)&\beta_{23}(\bx)\\
a_4+\beta_{31}(\bx)&\beta_{32}(\bx)&\beta_{33}(\bx)\end{pmatrix},
\end{equation}
where each $\beta_{kl}:\mathbb{R}^3\longrightarrow \mathbb{R}$ is a continuous function supported on $\mathbb{U}_2$ and satisfying  
$\sup\bigl\{|\beta_{kl}(\bx)|\, ;\, \bx \in \mathbb{R}^3\bigr\}=O(\varepsilon^2)$.
We set again $D\varphi_{g,2}(\bx)(v^{\rm u},v^{\rm s},v^{\rm cs})^T
=(\widehat{v}^{\rm u},\widehat{v}^{\rm s},\widehat{v}^{\rm cs})^T$.
As in the previous case, 
\begin{equation}\label{eqn_vs_a3}
|\widehat{v}^{\rm s}|\leq \varepsilon |a_3|+O(\varepsilon^2)
\end{equation}
for any $\bv=(v^{\mathrm{u}},v^{\mathrm{s}},v^{\mathrm{cs}})\in 
\boldsymbol{C}_{\varepsilon}^{\mathrm{uc}}(\boldsymbol{x})$ with 
$(v^{\mathrm{u}})^2+(v^{\mathrm{cs}})^2=1$.

First we consider the case when $|v^{\rm u}|\geq \dfrac{a_2}{2a_1}|v^{\rm cs}|$ or equivalently 
$|v^{\rm u}|\geq \dfrac{a_2}{\sqrt{4a_1^2+a_2^2}}$.
By \eqref{eqn_min} and \eqref{eqn_Dphi2}, it holds that 
\[
|\widehat{v}^{\rm cs}|\geq |a_4v^{\rm u}|-O(\varepsilon^2)\geq \dfrac{a_2|a_4|}{\sqrt{4a_1^2+a_2^2}}-O(\varepsilon^2)>|a_3|
\]
if $\varepsilon$ is taken sufficiently small.
By \eqref{eqn_vs_a3}, we may also assume that 
\begin{equation}\label{eqn_vcs>v3}
\varepsilon |\widehat{v}^{\rm cs}|>|\widehat{v}^{\rm s}|.
\end{equation}



Next we suppose that $|v^{\rm u}|<\dfrac{a_2}{2a_1}|v^{\rm cs}|$ or equivalently 
$|v^{\rm cs}|> \dfrac{2a_1}{\sqrt{4a_1^2+a_2^2}}$.
Again by \eqref{eqn_min} and \eqref{eqn_d_eta}, \eqref{eqn_Dphi2}, 
\begin{align*}
|\widehat{v}^{\rm u}|&\geq a_2|v^{\rm cs}|-\left|\frac{d\varphi}{dx}(x)v^{\rm u}\right|-O(\varepsilon^2)
\geq a_2|v^{\rm cs}|-a_2\left(\frac12+\varepsilon_0\right)|v^{\rm cs}|-O(\varepsilon^2)\\
&> a_2\left(\frac12-\varepsilon_0\right)\dfrac{2a_1}{\sqrt{4a_1^2+a_2^2}}-O(\varepsilon^2)
=\dfrac{a_1a_2}{\sqrt{4a_1^2+a_2^2}}(1-2\varepsilon_0)-O(\varepsilon^2)\\
&>|a_3|-O(\varepsilon^2).
\end{align*}
Thus we have 
\begin{equation}\label{eqn_vu>v3}
\varepsilon |\widehat{v}^{\rm u}|>|\widehat{v}^{\rm s}|
\end{equation}
for any sufficiently small $\varepsilon$. 
By \eqref{eqn_vcs>v3} and \eqref{eqn_vu>v3}, 
$|\widehat{v}^{\rm s}|\leq \varepsilon\sqrt{(\widehat{v}^{\rm u})^2+(\widehat{v}^{\rm cs})^2}$.
It follows that  
$D\varphi_{g,2}(\bx)\bv\in \boldsymbol{C}_{\varepsilon}^{\mathrm{uc}}(\boldsymbol{x})$.


Since $\varphi_{g,i}=\varphi_i$ on $\mathbb{R}^3\setminus \uu_i$ for $i=0,1,2$, 
the restriction 
$\mathrm{proj}|_{\varphi_{g,i}(S)}:\varphi_{g,i}(S)\longrightarrow \mathbb{R}^2$ is proper and surjective.
Hence $\varphi_{g,i}(S)$ is an $\varepsilon$-almost horizontal proper surface.
This completes the proof.
\end{proof}


\section{Proof of Theorem \ref{thm_A}}\label{S_proof_of_thmA}




A subset $A$ of $\mathbb{R}^3$ is $\varepsilon$-\emph{almost horizontal} if $A$ is contained in 
an $\varepsilon$-almost horizontal proper surface $S$ in $\mathbb{R}^3$.
If such a subset $A$ is Lebesgue measurable in $S$, then by Lemma \ref{l_ea_horizontal} 
\begin{equation}\label{eqn_Leb_A}
\mathrm{Leb}_S(A)\geq \mathrm{Leb}\,(\widehat A)\geq \frac{\mathrm{Leb}_S(A)}{\sqrt{1+\varepsilon^2}}
=\mathrm{Leb}_S(A)(1-O(\varepsilon^2)).
\end{equation}
Here $\mathrm{Leb}_S(\cdot)$ denotes the Lebesgue measure on $S$ with respect to 
the Riemannian metric on $S$ induced from the Euclidean metric on $\mathbb{R}^3$ 
and $\mathrm{Leb}\,(\cdot)$ does the standard Lebesgue measure on $\mathbb{R}^2$.



Recall that, for any $g\in \mathcal{U}_0$, the continuation $\Lambda_g$ of $\Lambda_{f_0}$ is associated with the coding pair $\{ \mathbb{U}_0, \mathbb{U}_1 \}$ given in Section \ref{S_near_f0}.
Let $\mathcal{I}_g: \Lambda_{g}  \longrightarrow  \{ 0,1\} ^{\mathbb Z}$ be the coding map 
determined from the pair. 
We say that an element $\bx_g$ of $\Lambda_g$ with $\mathcal{I}_g(\bx_g)=(v_j)_{j\in \mathbb{Z}}$ satisfies the \emph{asymptotically 1-filling condition} if 
\begin{equation}\label{eqn_1-filling}
\lim_{n\to \infty}\frac{\#\left\{j\,;\,0\leq j\leq n-1,\ v_j=1\right\}}{n}=1.
\end{equation}
We note that the subset of $\Lambda_g$ consisting of elements with the asymptotically 1-filling  
condition is dense in $\Lambda_g$.


\begin{remark}\label{r_ab+2}
As in the proof of \cite[Lemma 2.2]{KNS}, we may assume that $\alpha_k$ and $\beta_k$ in Definition \ref{dfn1} satisfy  
\begin{equation}\label{eqn_ab+2}
\alpha_k+\beta_k+2\leq \alpha_{k+1}
\end{equation}
for any $k\geq 1$.
In the case when $\alpha_k+\beta_k+1=\alpha_{k+1}$, we consider the new interval 
$[\alpha_k,\beta_{k+1}]\cap \mathbb{Z}$ instead of $\mathbb{I}_k\cup \mathbb{I}_{k+1}$. 
Repeating the construction, we have a sequence of intervals, still denoted by $(\mathbb{I}_k)$,  which satisfies not only \eqref{D1}, \eqref{D2} but also \eqref{eqn_ab+2}.
When $(\mathbb{I}_k)$ consists of only finitely many intervals, 
the last entry $\mathbb{I}_{k_0}$ is a half open interval $[\alpha_{k_0},\infty)\cap 
\mathbb{Z}$.
Then we split $\mathbb{I}_{k_0}$ into infinitely many intervals such that
$\mathbb{I}_{k_0}^{\mathrm{new}}=[\alpha_{k_0},\alpha_{k_0}+2]\cap \mathbb{Z}$ and 
$\mathbb{I}_{k_0+i}^{\mathrm{new}}=[\alpha_{k_0}+2^{i+1},\alpha_{k_0}
+2^{i+2}-2]\cap\mathbb{Z}$ for $i\geq 1$, 
which still satisfies \eqref{D1}, \eqref{D2} and \eqref{eqn_ab+2}.
\end{remark}




\begin{proof}[Proof of Theorem \ref{thm_A}]
The proof is done by contradiction.
Suppose that there exists an element $g_*$ of $\mathcal{U}_0$ which is strongly 
pluripotent for the continuation $\Lambda_{g_*}$ of the whole $\Lambda_{f_0}$.
Take an element $\bx_*\in \Lambda_{g_*}$ with asymptotically 1-filling condition.
We set $\mathcal{I}_{g_*}(\bx_*)=(v_j)_{j\in \mathbb{Z}}$ and $\Sigma=\{v\}$ 
for $v=(v_j)_{j\in \mathbb{N}_0}$.
Since $g_*$ is strongly pluripotent for the subset $\mathcal{I}_{g_*}^{-1}(\widehat\Sigma)$ of the whole $\Lambda_{g_*}$, 
$g_*$ is $\Sigma$-describable by Theorem \ref{p-lemma}, that is,  
there exist $g\in \mathcal{U}_0$ arbitrarily $C^r$-close to $g_*$,  
a subset $D_g$ of $M$ with positive Lebesgue measure and 
non-negative integers $\alpha_k$, $\beta_k$ satisfying \eqref{D1}, \eqref{D2} 
and \eqref{eqn_ab+2}.


By using \eqref{D1}, one can show that $\beta_{k}\neq 0$ for infinitely many $k\in \mathbb{N}$.
In fact, if $\beta_{k}=0$ for all but finitely many $k\in \mathbb{N}$, then by \eqref{eqn_ab+2} 
we would have 
\[
\limsup _{n\rightarrow \infty} \frac{\#\left\{j\,;\, 0\le j \le n-1,\ j \in \bigcup _{k=1}^\infty \mathbb{I}_k\right\}}{n}\leq \frac12,
\]
which contradicts \eqref{D1}.
In particular we may assume that $\beta_1\geq 1$ by regarding $k-k_0+1$ as a new $k$ for some $k_0$ with 
$\beta_{k_0}\neq 0$.
Since $g^{\alpha_1}(D_g)\subset \uu_a$ and $g^{\alpha_1+1}(D_g)\subset \uu_b$ for some $a,b\in \{0,1\}$, 
$g^{\alpha_1+1}(D_g)$ is contained in $g(\uu_a)\cap \uu_b\subset \mathbb{B}$, 
where $\mathbb{U}_i$ $(i=0,1)$ is the open neighborhood of $\xx_{g,i}$ in $M$ given in Section \ref{S_near_f0}.
Let $\mathbb{T}_g$ be the union $\xx_{g,0}\cup \xx_{g,1}\cup \xx_{g,2}\cup g(\xx_{g,2})$. 
For any $j\geq \alpha_1+1$, we will see that $g^j(D_g)$ is a subset of $\mathbb{T}_g$.
If not, there would exist $j\geq \alpha_1+1$ with $g^j(D_g)\setminus \mathbb{T}_g\neq \emptyset$.
We denote by $j_0$ the smallest integer among such $j$. 
Then $g^{j_0}(D_g)$ contains an element $\by$ of $\yy_g\cup \mathbb{S}_g\setminus 
\mathrm{Fr}(\yy_g\cup \mathbb{S}_g)$, 
which is obvious when $j_0=\alpha_1+1$ and proved by using the fact $g^{j_0-1}(D_g)\subset \mathbb{T}_g$ 
when $j_0>\alpha_1+1$.
By Extra Assumptions with respect to $g$, $(g^{j_0+n}(\by))_{n\geq 0}$ converges to the attracting fixed point $s_{g,0}$ of $g$ which is the continuation of $s_0$.
On the other hand, by \eqref{D2}, $g^{\alpha_k-j_0}(\by)\in g^{\alpha_k}(D_g)$ 
is contained in either $\mathbb{U}_{0}$ or $\mathbb{U}_{1}$ for any $\alpha_k>j_0$, a contradiction.
Thus we have $g^j(D_g)\subset \mathbb{T}_g$ for any $j\geq \alpha_1+1$.
This implies that $g^{\alpha_k}(D_g)$ is in either $\xx_{g.0}$ or $\xx_{g,1}$ for any $k\geq 2$.





By Fubini's Theorem, there exists $y_0\in (-\varepsilon_0,1+\varepsilon_0)$ such that 
the horizontal section $A_{\alpha_2}=g^{\alpha_2}(D_g)\cap \{(x,y,z)\in\mathbb{B}\,;\, y=y_0\}$ has positive 2-dimensional Lebesgue measure.


We set
\[\gamma_k=\alpha_{k+1}-(\alpha_k+\beta_k)\]
and define inductively subsets $A_j$ $(j> \alpha_2)$ of $M$ each of which is contained in a component of $\mathbb{T}_g$.

First suppose that $A_j$ $(\alpha_2\leq j\leq \alpha_k)$ are already defined for some $k\geq 2$.
\begin{enumerate}[(1)]
\item
Let $A_j=g^{j-\alpha_k}(A_{\alpha_k})$ for $j=\alpha_k+1,\dots,\alpha_k+\beta_k$, 
which are contained in either $\mathbb{X}_{g,0}$ or $\mathbb{X}_{g,1}$.
\item
If $A_{\alpha_k+\beta_k+j-1}$ is contained in $\xx_{g,i}$ $(i=0,1)$ 
for some $j$ with $1\leq j< \gamma_k$, then 
$g(A_{\alpha_k+\beta_k+j-1})$ is a disjoint  union of $C_{\alpha_k+\beta_k+j;i}=
g(A_{\alpha_k+\beta_k+j-1})\cap \mathbb{X}_{g,i}$ $(i=0,1,2)$.
Possibly some of them are empty.
We suppose that $A_{\alpha_k+\beta_k+j}$ is one of $C_{\alpha_k+\beta_k+j;i}$ with 
$\mathrm{Leb}\,(\widehat C_{\alpha_k+\beta_k+j;i})\geq \mathrm{Leb}\,(\widehat C_{\alpha_k+\beta_k+j;s})
$ 
for any $s\in \{0,1,2\}\setminus \{i\}$.
In particular, 
\[
\mathrm{Leb}\,(\widehat A_{\alpha_k+\beta_k+j})\geq \frac13\mathrm{Leb}\bigl(\widehat{g(A_{\alpha_k+\beta_k+j-1})}\bigr).
\]
\item
If 
$A_{\alpha_k+\beta_k+j-1}$ is contained in $\xx_{g,2}$ for some $j$ with $1\leq j< \gamma_k-1$, then we set 
$A_{\alpha_k+\beta_k+j}=g(A_{\alpha_k+\beta_k+j-1})$ and $A_{\alpha_k+\beta_k+j+1}=g^2(A_{\alpha_k+\beta_k+j-1})$.
\end{enumerate}
Let $\gamma_k^{(3)}$ denote the number of integers $j$ with $1 \le j < \gamma_k - 1$ for which (3) holds.

By Lemma \ref{l_ea_horizontal}, any $A_j$ in $\mathbb{T}_g\setminus g(\xx_{g,2})$ is $\varepsilon$-almost horizontal.
By \eqref{eqn_Leb_P0}, \eqref{eqn_Leb_P} and \eqref{eqn_Leb_A}, there exist constants 
\[
\rho_i=\lambda_{{\rm cs}i}\lambda_{\rm u}-O(\varepsilon^2)\quad (i=0,1),\quad 
\rho_2=|a_2a_4|-O(\varepsilon^2)
\]
such that 
\begin{equation}\label{eqn_LebA}
\begin{split}
\mathrm{Leb}\,(\widehat A_{j+1})&\geq \rho_i\mathrm{Leb}\,(\widehat A_{j})\\
&\qquad\text{if $\alpha_k\leq j\leq \alpha_k+\beta_k-1$ and $A_{j}\subset \mathbb{X}_{g,i}$ $(i=0,1)$},\\
\mathrm{Leb}\,(\widehat A_{j+1})&\geq \frac13\rho_i\mathrm{Leb}\,(\widehat A_{j})\\
&\qquad\text{if $\alpha_k+\beta_k\leq j\leq \alpha_{k+1}-1$ and $A_{j}\subset \mathbb{X}_{g,i}$ $(i=0,1)$},\\
\mathrm{Leb}\,(\widehat A_{j+2})&\geq \rho_2\mathrm{Leb}\,(\widehat A_{j})\\
&\qquad\text{if $\alpha_k+\beta_k+1\leq j\leq \alpha_{k+1}-2$ and $A_{j}\subset \mathbb{X}_{g,2}$.}
\end{split}
\end{equation}
Here we note that $A_{\alpha_{k+1}-1}$ is not contained in $\xx_{g,2}$.
If not, by Extra Assumptions $A_{\alpha_{k+1}}\subset g(\xx_{g,2})$ would be disjoint from $\xx_{g,0}\cup \xx_{g,1}$, a contradiction.
By \eqref{eqn_lamcs1}, one can suppose that $\rho_1>1$.
Now we consider the constant $\widehat\rho=\min\left\{\dfrac{\rho_0}3,\dfrac{\rho_1}3,\rho_2,1\right\}$.
By Definition \ref{describable}\,\eqref{D1}, $\lim\limits_{k\to \infty}\sum_{i=2}^k \gamma_i/\sum_{i=2}^k\beta_i=0$.
Let $\tau_k=\#\{j\in \mathbb{I}_k\,;\,v_j=0\}$.
Then the asymptotically 1-filling condition \eqref{eqn_1-filling} implies $\lim\limits_{k\to \infty}\sum_{i=2}^k \tau_i/\sum_{i=2}^k\beta_i=0$ 
and hence in particular $\frac12\sum_{i=2}^k \beta_i\geq \sum_{i=2}^k\tau_i$ for all sufficiently large $k$.
Since $\rho_1>1$ and $\widehat\rho\leq 1$, we have by \eqref{eqn_LebA} 
\begin{align*}
\mathrm{Leb}\,(\widehat A_{\alpha_{k+1}})&\geq \rho_1^{\,\sum_{i=2}^k(\beta_i-\tau_i)}
\widehat\rho^{\,\sum_{i=2}^k(\gamma_i-\gamma_i^{(3)}+\tau_i)}\mathrm{Leb}\,(A_{\alpha_2})\\
&\geq \rho_1^{\,\frac12\sum_{i=2}^k\beta_i}
\widehat\rho^{\,\sum_{i=2}^k(\gamma_i+\tau_i)}\mathrm{Leb}\,(A_{\alpha_2}).
\end{align*}
Here `$-\gamma_i^{(3)}$' arises from skipping the step 
$A_{\alpha_k+\beta_k+j} = g(A_{\alpha_k+\beta_k+j-1})$ in each case of (3).
It follows that 
\begin{align*}
\log \mathrm{Leb}\,(\widehat A_{\alpha_{k+1}})&\geq 
%\sum_{i=2}^k(\beta_i-\tau_i)\log \rho_1+
%\sum_{i=2}^k(\gamma_i+\tau_i)\log \widehat\rho+\log\mathrm{Leb}\,(A_{\alpha_2})\\
%&\geq 
\frac12\sum_{i=2}^k\beta_i\log \rho_1+
\sum_{i=2}^k(\gamma_i+\tau_i)\log \widehat\rho+\log\mathrm{Leb}\,(A_{\alpha_2})\\
&=\frac12\sum_{i=2}^k\beta_i\left(\log \rho_1+\frac{2\sum_{i=2}^k(\gamma_i+\tau_i)}{\sum_{i=2}^k\beta_i}\log \widehat\rho\right)+\log\mathrm{Leb}\,(A_{\alpha_2})\to \infty
\end{align*}
as $k\to \infty$ and so $\lim\limits_{k\to \infty}\mathrm{Leb}\,(\widehat A_{\alpha_{k+1}})=\infty$, which contradicts that $\widehat A_{\alpha_{k+1}}$ is a subset of $I_{\varepsilon_0}^2$.
This completes the proof.
\end{proof}


\begin{proof}[Proof of Corollary \ref{cor_B}]
Let $f_0$ be a diffeomorphism satisfying all the conditions as above 
for the proof of Theorem \ref{thm_A} and such that $\Lambda_{f_0}$ is a wild blender-horseshoe.
Moreover we suppose that $f_0$ employs positive constants $\lambda_{\rm ss}, \lambda_{\rm cs0}, \lambda_{\rm cs1}$ and $\lambda_{\rm u}$ 
satisfying 
\begin{equation}\label{eqn_KNS}
\lambda_{\rm ss}<\lambda_{\rm cs0}<\frac12<\lambda_{\rm cs1}<1<
\lambda_{\rm cs0}+\lambda_{\rm cs1},\quad 2<\lambda_{\rm u},\quad \lambda_{\rm cs0}\lambda_{\rm cs1}\lambda_{\rm u}^2<1
\end{equation}
instead of \eqref{eqn_eigen_v}, which are 
the same as the conditions (3.1a) and (3.1b) in \cite{KNS}.
We recall that Theorem B in \cite{KNS} works under $C^r$-category with $2\leq r<\infty$, 
which is an assumption of Corollary \ref{cor_B}.
Hence any element $f$ of $\mathcal{U}_0$ is strongly pluripotent 
for $\Lambda_f^{({\rm mj})}$.
Since \eqref{eqn_KNS} implies \eqref{eqn_eigen_v},  
$f$ is not strongly pluripotent for $\Lambda_f$.
\end{proof}



\section*{Acknowledgement}
S.~Kiriki was supported by JSPS KAKENHI Grant Number 21K03332.
X.~Li was supported by NSFC Grant Numbers 11701199 and 12331005.
Y.~Nakano was supported by JSPS KAKENHI Grant  Number 23K03188 and JST PRESTO Grant Number JPMJPR25K8.
T.~Soma was supported by JSPS KAKENHI Grant Number 22K03342.



\begin{thebibliography}{9}



\bibitem{B22}
P.~G. Barrientos, 
\newblock Historic wandering domains near cycles, 
\newblock {\em Nonlinearity}, 35(6):3191--3208, 2022.

\bibitem{BB23}
P.~Berger and S.~Biebler, 
\newblock Emergence of wandering stable components, 
\newblock {\em J. Amer. Math. Soc.}, 36(2):397--482, 2023.

\bibitem{BD96}
Ch. Bonatti and L.~J. D{\'{\i}}az, 
\newblock Persistent nonhyperbolic transitive diffeomorphisms, 
\newblock {\em Ann. of Math. (2)}, 143(2):357--396, 1996.

\bibitem{BD12}
Ch. Bonatti and L.~J. D{\'{\i}}az, 
\newblock Abundance of {$C^1$}-robust homoclinic tangencies, 
\newblock {\em Trans. Amer. Math. Soc.}, 364(10):5111--5148, 2012.

\bibitem{BDV05}
Ch. Bonatti, L.~J. D{\'{\i}}az, and M.~Viana, 
\newblock {\em Dynamics beyond uniform hyperbolicity}, volume 102 of {\em
  Encyclopaedia of Mathematical Sciences}, 
\newblock Springer-Verlag, Berlin, 2005, 
\newblock A global geometric and probabilistic perspective, Mathematical
  Physics, III.

\bibitem{CV01}
E.~Colli and E.~Vargas, 
\newblock Non-trivial wandering domains and homoclinic bifurcations, 
\newblock {\em Ergodic Theory Dynam. Systems}, 21(6):1657--1681, 2001.

\bibitem{JM}
Uygun Jamilov and Farrukh Mukhamedov, 
\newblock Takens' the last problem and {S}tein-{U}lam spiral type maps, 
\newblock {\em J. Math. Anal. Appl.}, 531(1):Paper No. 127813, 13, 2024.

\bibitem{KLNSV}
Shin Kiriki, Xiaolong Li, Yushi Nakano, Teruhiko Soma, and Edson Vargas, 
\newblock Takens' {L}ast {P}roblem and strong pluripotency, 
\newblock {\em Nonlinearity}, 38(8):Paper No. 085007, 2025.

\bibitem{KLNS22}
Shin Kiriki, Xiaolong Li, Yushi Nakano, and Teruhiko Soma, 
\newblock Abundance of observable {L}yapunov irregular sets, 
\newblock {\em Comm. Math. Phys.}, 391(3):1241--1269, 2022.

\bibitem{KNS}
S.~Kiriki, Y.~Nakano, and T.~Soma, 
\newblock Pluripotency of wandering dynamics, 
\newblock {\em arXiv2404.00337}.

\bibitem{KNS23}
S.~Kiriki, Y.~Nakano, and T.~Soma, 
\newblock Historic and physical wandering domains for wild blender-horseshoes, 
\newblock {\em Nonlinearity}, 36(8):4007--4033, 2023.

\bibitem{KS17}
S.~Kiriki and T.~Soma, 
\newblock Takens' last problem and existence of non-trivial wandering domains.
\newblock {\em Adv. Math.}, 306:524--588, 2017, 

\bibitem{LR17}
I.~S. Labouriau and A.~A.~P. Rodrigues, 
\newblock On {T}akens' last problem: tangencies and time averages near
  heteroclinic networks, 
\newblock {\em Nonlinearity}, 30(5):1876--1910, 2017.

\bibitem{R01}
D.~Ruelle, 
\newblock {\em Historical behaviour in smooth dynamical systems}, pages 63--66, 
\newblock Inst. Phys., Bristol, 2001.

\bibitem{T08}
F.~Takens, 
\newblock Orbits with historic behaviour, or non-existence of averages, 
\newblock {\em Nonlinearity}, 21(3):T33--T36, 2008.




\end{thebibliography}












\end{document}
